\documentclass[10pt,a4paper]{article}

\usepackage[margin=22mm]{geometry}
\usepackage{amsmath,amssymb,mathtools}
\usepackage{microtype}
\usepackage[hidelinks]{hyperref}
\usepackage{fancyhdr}

\setlength{\parindent}{0pt}
\setlength{\parskip}{3pt}
\setlength{\abovedisplayskip}{6pt}
\setlength{\belowdisplayskip}{6pt}
\setlength{\abovedisplayshortskip}{3pt}
\setlength{\belowdisplayshortskip}{3pt}

\pagestyle{fancy}
\fancyhf{}
\fancyhead[L]{Free-space LFM receive beamforming}
\fancyhead[R]{Vertical line array}
\fancyfoot[C]{\thepage}
\renewcommand{\headrulewidth}{0.4pt}

\newcommand{\jj}{\mathrm{j}}
\newcommand{\e}{\mathrm{e}}
\newcommand{\dd}{\mathrm{d}}

\title{\vspace{-1.3em}\textbf{Far-Field LFM Receive Beamforming in Free Space}\\
\large Point source and vertical line array}
\author{}
\date{}

\begin{document}
\maketitle
\vspace{-2.5em}

This note describes conventional broadband receive beamforming for one point source
and a vertical line array (VLA) in homogeneous free space. Depth is positive downward,
and the look angle is measured from the horizontal with positive angles pointing
downward. True time delays are used across the full LFM bandwidth.

\section{LFM waveform and far-field geometry}

For start frequency $f_0$, stop frequency $f_1$, and duration $T$, the transmitted
real LFM waveform is
\begin{equation}
 s(t)=g(t)\cos\!\left[2\pi\left(f_0t+\frac{f_1-f_0}{2T}t^2\right)\right],
 \qquad 0\leq t<T,
\end{equation}
where $g(t)$ is a short edge taper.

Let receiver $n$ be at depth $z_n$. Define the VLA reference depth and relative
element coordinate by
\begin{equation}
 z_0=\frac{1}{N}\sum_{n=1}^{N}z_n,
 \qquad \widetilde z_n=z_n-z_0.
\end{equation}
For a source at horizontal range $r_s$ and depth $z_s$, its range and bearing at the
array reference are
\begin{equation}
 R=\sqrt{r_s^2+(z_s-z_0)^2},
 \qquad
 \theta_s=\operatorname{atan2}(z_s-z_0,r_s).
\end{equation}
The exact and far-field source-to-element ranges are
\begin{equation}
 R_n=\sqrt{r_s^2+(z_s-z_n)^2}
 \simeq R-\widetilde z_n\sin\theta_s.
 \label{eq:far-field-range}
\end{equation}
The first neglected term is
$\widetilde z_n^2\cos^2\theta_s/(2R)$, which provides a direct check on the
plane-wave approximation.

Define the relative arrival delay for a candidate look angle $\theta$ as
\begin{equation}
 d_n(\theta)=-\frac{\widetilde z_n\sin\theta}{c}.
 \label{eq:relative-delay}
\end{equation}
Using Eq.~\eqref{eq:far-field-range}, the signal at receiver $n$ is
\begin{equation}
 x_n(t)=\beta s\!\left(t-\frac{R}{c}-d_n(\theta_s)\right)+v_n(t),
 \qquad \beta=\frac{A}{R},
 \label{eq:received-signal}
\end{equation}
where $A$ is the source-strength constant and $v_n(t)$ is receiver noise. Spherical
spreading is retained through $1/R$, while its small variation across the aperture is
neglected in the far field.

\section{True-time-delay beamformer}

For normalized spatial weights $q_n$ satisfying $\sum_nq_n=1$, the ideal offline
beamformer steered toward $\theta_0$ is
\begin{equation}
 y(t;\theta_0)=\sum_{n=1}^{N}q_n
 x_n\!\left(t+d_n(\theta_0)\right).
 \label{eq:noncausal-beamformer}
\end{equation}
At $\theta_0=\theta_s$, every relative delay in Eq.~\eqref{eq:received-signal}
cancels and the LFM signals add coherently.

A causal implementation uses only nonnegative delays:
\begin{equation}
 d_{\max}(\theta_0)=\max_n d_n(\theta_0),
 \qquad
 D_n(\theta_0)=d_{\max}(\theta_0)-d_n(\theta_0),
\end{equation}
\begin{equation}
 y_c(t;\theta_0)=\sum_{n=1}^{N}q_nx_n(t-D_n).
 \label{eq:causal-beamformer}
\end{equation}
Equation~\eqref{eq:causal-beamformer} differs from
Eq.~\eqref{eq:noncausal-beamformer} only by the common delay $d_{\max}$. Fractional
values of $D_n$ require a fixed fractional-delay filter, such as Lagrange, Farrow,
windowed-sinc, or a properly padded FFT implementation.

With Fourier convention $X(f)=\int x(t)\e^{-\jj2\pi ft}\,\dd t$, the equivalent
frequency-domain beamformer is
\begin{equation}
 Y(f;\theta_0)=\sum_{n=1}^{N}q_nX_n(f)
 \e^{\jj2\pi f d_n(\theta_0)}.
 \label{eq:frequency-beamformer}
\end{equation}
The linear phase ramp in Eq.~\eqref{eq:frequency-beamformer} is a true time delay.
A single phase shift evaluated at the center frequency would cause beam squint.
Equivalently, with
$[\boldsymbol a(f,\theta)]_n=\e^{-\jj2\pi f d_n(\theta)}$, the conventional output
is $Y(f;\theta)=\boldsymbol a^{\mathrm H}(f,\theta)\boldsymbol X(f)/N$ for uniform
weights $q_n=1/N$.

\section{Broadband scan, array response, and matched filtering}

For FFT frequencies $f_k$ in the LFM band $\mathcal B$, the implemented broadband
energy scan is
\begin{equation}
 P_E(\theta)=\sum_{f_k\in\mathcal B}
 \left|\sum_{n=1}^{N}q_nX_n(f_k)
 \e^{\jj2\pi f_kd_n(\theta)}\right|^2.
 \label{eq:broadband-scan}
\end{equation}
This is an incoherent sum over frequency of coherently beamformed bin powers. The
displayed scan is
\begin{equation}
 P_{\mathrm{dB}}(\theta)=10\log_{10}
 \frac{P_E(\theta)}{\max_{\vartheta}P_E(\vartheta)}.
\end{equation}

For a plane wave from $\theta$ while steering toward $\theta_0$, the ideal response is
\begin{equation}
 B(f,\theta;\theta_0)=\sum_{n=1}^{N}q_n
 \e^{\jj2\pi f\widetilde z_n(\sin\theta-\sin\theta_0)/c}.
\end{equation}
True-time-delay steering keeps the maximum at $\theta_0$ throughout the band, but
the main-lobe width decreases with frequency.

For known waveform $s(t)$, the pulse-compressed output is
\begin{equation}
 z(\tau;\theta_0)=\int y_c(t;\theta_0)s^*(t-\tau)\,\dd t,
 \qquad
 Z(f;\theta_0)=Y_c(f;\theta_0)S^*(f).
\end{equation}
Beamforming supplies spatial gain; matched filtering supplies the separate LFM
time-bandwidth gain. Since both are linear, their order can be interchanged when all
channels use the same temporal filter.

For independent equal-variance receiver noise and $q_n=1/N$, the coherent signal
amplitude is unchanged while the output noise variance becomes $\sigma_v^2/N$.
Thus
\begin{equation}
 \frac{\mathrm{SNR}_{\mathrm{out}}}{\mathrm{SNR}_{\mathrm{in}}}=N,
 \qquad G_{\mathrm{array}}=10\log_{10}N\ \mathrm{dB}.
\end{equation}
For the $N=16$ default simulation, the ideal spatial gain is $12.04\ \mathrm{dB}$.

Finally, full-sector unambiguous steering and conservative far-field operation require
\begin{equation}
 \Delta z\leq\frac{c}{2f_{\max}},
 \qquad
 R\gtrsim\frac{2D^2}{\lambda_{\min}},
 \qquad \lambda_{\min}=\frac{c}{f_{\max}},
\end{equation}
where $D$ is the VLA aperture. A VLA measures only the vertical direction cosine;
azimuth remains unresolved in free space.

\end{document}
